\documentclass[12pt]{article}

\input{short-template.tex}

\newcommand{\R}{\mathbb{R}}
\newcommand{\N}{\mathcal{N}}

\begin{document}

\MakeShortHeader
  {LeWorldModel}
  {SIGReg 为什么能防塌缩？}
  {World Model,SIGReg,Latent Planning,JEPA}

\begin{KeyBox}
世界模型的基本想法是：\Emph{根据当前的状态，采取一个动作之后，能够预测世界下一个时刻的状态}。  
LeWM 属于 JEPA 路线：不重建像素，而是直接预测下一步 latent。
\end{KeyBox}

\section*{世界模型}

世界模型根据其基本想法，衍生出了许多不同的路线，部分人会把预测视频的下一帧也并入世界模型的范畴。  
而JEPA路线则是先把高维观测压到 latent 空间，再在 latent 里学习 dynamics。这样做的好处是：预测和规划都发生在更紧凑的表示空间里，计算更轻，也更接近控制本身关心的结构。但是与视频相比坏处就是无法直接刻画表征的好坏。

LeWM是从原始像素出发，端到端学习一个\Emph{可预测、可规划、但不塌缩}的 latent world model。

\section*{LeWM}

\begin{FormulaBox}[训练目标]
\[
    z_t=\operatorname{enc}_\theta(o_t),
    \qquad
    \hat z_{t+1}=\operatorname{pred}_\phi(z_t,a_t).
\]
\[
    \mathcal L_{\mathrm{pred}}
    = \frac{1}{T-1}\sum_{t=1}^{T-1}
      \lVert \hat z_{t+1}-z_{t+1}\rVert_2^2,
    \qquad
    \mathcal L_{\mathrm{LeWM}}
    = \mathcal L_{\mathrm{pred}}+\lambda\,\mathrm{SIGReg}(Z).
\]
\end{FormulaBox}
 
$z_t$就可以当成$t$时刻的“状态”，$o_t$是我们观测到的视觉信息（像素），$a_t$是采取的动作。

\section*{防止塌缩}

假设 encoder 把所有观测都映到同一个向量 $c$，也就是
\[
    z_t\equiv c,
    \qquad
    \hat z_{t+1}\equiv c.
\]
那么在欧氏空间下，预测损失立刻变成
\[
    \lVert \hat z_{t+1}-z_{t+1}\rVert_2^2=0.
\]

这说明：\Emph{在 next-latent prediction 这个目标下，点塌缩不是“训练失败”，而是一个合法的最优解。}  
一旦如此，所有状态在 latent 空间中都不可区分。当前状态、目标状态、预测终点都会落在同一个点上，于是任何动作序列看起来都一样好，规划也就失去了意义。

\section*{SIGReg}

记 latent 样本矩阵为
\[
    Z\in\R^{n\times d}.
\]
SIGReg 的目标是更强的一件事：\Emph{把 latent 的经验分布推向各向同性高斯 $\N(0,I_d)$}。

\begin{FormulaBox}[SIGReg]
采样随机单位方向
\[
    u^{(m)}\sim \operatorname{Unif}(\mathbb{S}^{d-1}),
    \qquad m=1,\dots,M,
\]
做一维投影
\[
    h^{(m)}=Zu^{(m)}\in\R^n,
\]
然后平均一维正态性统计量：
\[
    \mathrm{SIGReg}(Z)=\frac{1}{M}\sum_{m=1}^M T\!\left(h^{(m)}\right).
\]
其中 $T$ 是 Epps--Pulley 统计量，用来衡量投影样本与标准正态 $\N(0,1)$ 的距离。
\end{FormulaBox}

理论依据是：\Emph{一个高维分布，可以由它在所有一维方向上的投影分布刻画。}  
所以，高维 anti-collapse 问题可以转化成很多个一维问题：随机看几个方向，检查投影后是否是近似标准高斯。这样一来，常值解、低秩表示、强相关表示都会被排除。

\section*{Latent Planning}

训练结束后，固定 encoder 和 predictor。  
给定当前观测 $o_1$ 与目标观测 $o_g$，先编码为
\[
    \hat z_1=\operatorname{enc}_\theta(o_1),
    \qquad
    z_g=\operatorname{enc}_\theta(o_g).
\]
然后在 latent 空间里滚动预测：
\[
    \hat z_{t+1}=\operatorname{pred}_\phi(\hat z_t,a_t),
    \qquad t=1,\dots,H-1.
\]

规划目标是
\begin{FormulaBox}[有限时域目标]
\[
    a_{1:H}^*
    =
    \arg\min_{a_{1:H}}
    \lVert \hat z_H-z_g\rVert_2^2.
\]
\end{FormulaBox}

\Term{planner 不修改 latent}。  
它真正优化的是一串动作，使 world model 预测出来的终点 latent，尽可能靠近目标观测的 latent 编码。

实现上使用的是 CEM 配合 MPC：在动作序列空间里反复采样、筛选、更新，再持续重规划，以减轻长时滚动造成的模型误差积累。

\begin{Takeaway}  
LeWM 先用预测损失学出动作条件下的 latent dynamics，再用 SIGReg 给 latent 空间施加非塌缩的分布几何；而 latent planning 则是在这套冻结的 dynamics 上，搜索一串动作，让预测终点尽量到达目标 latent。
\end{Takeaway}

\section*{参考文献}

\begin{enumerate}
    \item Lucas Maes, Quentin Le Lidec, Damien Scieur, Yann LeCun, Randall Balestriero. \textit{LeWorldModel: Stable End-to-End Joint-Embedding Predictive Architecture from Pixels}. arXiv preprint, 2026.

    \item David Ha, Jürgen Schmidhuber. \textit{World Models}. arXiv preprint arXiv:1803.10122, 2018.

    \item Yann LeCun. \textit{A Path Towards Autonomous Machine Intelligence}. Open Review, 2022.

    \item Randall Balestriero, Yann LeCun. \textit{LeJEPA: Provable and Scalable Self-Supervised Learning without the Heuristics}. arXiv preprint, 2025.

    \item Thomas W. Epps, Lawrence B. Pulley. \textit{A Test for Normality Based on the Empirical Characteristic Function}. Biometrika, 70(3):723--726, 1983.

    \item Harald Cram\'er, Herman Wold. \textit{Some Theorems on Distribution Functions}. Journal of the London Mathematical Society, 1(4):290--294, 1936.

    \item Reuven Y. Rubinstein, Dirk P. Kroese. \textit{The Cross-Entropy Method: A Unified Approach to Combinatorial Optimization, Monte-Carlo Simulation and Machine Learning}. Springer, 2004.
\end{enumerate}

\end{document}
